\section{Reproducibility}
\label{supp:reproducibility}

The build is organized around reproducible artifacts. The root command \texttt{make paper-assets} regenerates figures, tables, result macros, the results synthesis document, and per-file provenance records. The command \texttt{make paper} builds the main manuscript after regenerating assets. The command \texttt{make supplement} builds this supplement after the same asset step. The command \texttt{make paper-check} compiles both documents and audits generated assets, citations, compiled PDFs, and source markers. The command \texttt{make submission-package} assembles the reviewed manuscript package, reviewer report, revision matrix, availability statements, protocol, lock file, and artifact manifests without copying raw third-party datasets.

The generated macros provide stable result values used by the manuscript and supplement. Exact build identifiers and file-integrity records are stored in the repository manifests rather than printed in the narrative text.

Machine-readable interpretation and citation audits are written to \texttt{paper/generated/} by the paper checker. Manual edits to generated figures or tables are detected through the artifact-manifest audit.
